\documentclass[11pt]{article}
\usepackage[margin=1in]{geometry}
\usepackage[T1]{fontenc}
\usepackage[utf8]{inputenc}
\usepackage{amsmath,amssymb,amsthm}
\usepackage{enumitem}

\title{ICPC World Finals 2021\\H. Prehistoric Programs}
\author{}
\date{}

\begin{document}
\maketitle

\section*{Problem Summary}

Each tablet is a parentheses string. We may permute the tablets, then concatenate all strings. We need
the final concatenation to be a correct bracket sequence:
\begin{itemize}[leftmargin=*]
    \item every prefix must have nonnegative balance;
    \item the total balance must be zero.
\end{itemize}

\section*{Per-Tablet Statistics}

For one string, define:
\begin{itemize}[leftmargin=*]
    \item $\text{bal}$: total balance, where \texttt{(} is $+1$ and \texttt{)} is $-1$;
    \item $\text{minPref}$: the minimum prefix balance inside that string.
\end{itemize}

If the current balance before appending this string is $B$, then the string is safe to append exactly when
\[
    B + \text{minPref} \ge 0.
\]
After appending it, the new balance becomes
\[
    B + \text{bal}.
\]

\section*{Split into Two Groups}

Strings with nonnegative total balance are helpful: once we append them, the available balance never gets
smaller in the long run. Strings with negative total balance are dangerous and should be postponed.

So we split all tablets into:
\begin{itemize}[leftmargin=*]
    \item \textbf{positive group}: $\text{bal} \ge 0$;
    \item \textbf{negative group}: $\text{bal} < 0$.
\end{itemize}

If a valid order exists at all, there is one where all positive-group strings come first and all
negative-group strings come last.

\section*{How to Sort the Positive Group}

Inside the positive group, we want the least demanding strings first: the ones that require the smallest
starting balance. That means sorting by
\[
    \text{minPref}
\]
in decreasing order.

Then we greedily append them in that order, checking
\[
    B + \text{minPref} \ge 0.
\]

\section*{How to Sort the Negative Group}

For negative-balance strings, the right measure is not just $\text{minPref}$, because after appending the
string we lose some balance permanently.

The standard way to view this group is to reverse each string and swap open/close parentheses. This turns
it into a positive-balance problem. Under that transformation, the relevant key becomes
\[
    \text{bal} - \text{minPref}.
\]

Therefore the negative group should be sorted by
\[
    \text{bal} - \text{minPref}
\]
in decreasing order, and then greedily appended in that order.

\section*{Algorithm}

\begin{enumerate}[leftmargin=*]
    \item For each string, compute $\text{bal}$ and $\text{minPref}$.
    \item If the total balance over all strings is not zero, output \texttt{impossible}.
    \item Split strings into the positive and negative groups.
    \item Sort:
    \begin{itemize}[leftmargin=*]
        \item positive group by decreasing $\text{minPref}$;
        \item negative group by decreasing $\text{bal} - \text{minPref}$.
    \end{itemize}
    \item Traverse first the sorted positive group, then the sorted negative group.
    \item Maintain the current balance $B$. If some string has
    \[
        B + \text{minPref} < 0,
    \]
    output \texttt{impossible}.
    \item Otherwise append it, update
    \[
        B \leftarrow B + \text{bal},
    \]
    and record the index.
\end{enumerate}

\section*{Correctness Proof}

We prove that the algorithm outputs a valid permutation exactly when one exists.

\paragraph{Lemma 1.}
If a string with statistics $(\text{bal}, \text{minPref})$ is appended when the current balance is $B$, then
the concatenation remains prefix-valid if and only if
\[
    B + \text{minPref} \ge 0.
\]

\paragraph{Proof.}
Inside the string, every prefix changes the balance by exactly the corresponding prefix balance of that
string. The worst such drop is $\text{minPref}$. So all resulting prefixes stay nonnegative exactly when
the minimum possible one, namely $B+\text{minPref}$, is nonnegative. \qed

\paragraph{Lemma 2.}
If a valid ordering exists, then there also exists a valid ordering in which every string with
$\text{bal} \ge 0$ appears before every string with $\text{bal} < 0$.

\paragraph{Proof.}
A nonnegative-balance string never reduces the total available balance after it is appended, while a
negative-balance string does. Swapping a negative string that appears before a nonnegative string cannot
make the future balance smaller at the point where the nonnegative string starts. Repeating this exchange
process pushes all nonnegative strings to the front without destroying validity. \qed

\paragraph{Lemma 3.}
Among strings with $\text{bal} \ge 0$, sorting by decreasing $\text{minPref}$ is optimal for greedy
construction.

\paragraph{Proof.}
Take two such strings $A$ and $B$ with $\text{minPref}(A) \ge \text{minPref}(B)$. If some current balance
$X$ can append $B$ first, then
\[
    X + \text{minPref}(B) \ge 0.
\]
Since $\text{minPref}(A)$ is at least as large, $A$ is also appendable from $X$. After appending $A$, the
balance does not decrease overall because $\text{bal}(A)\ge0$, so whatever was possible after $B$ remains
possible after $A$. Thus placing the larger $\text{minPref}$ first is never worse. \qed

\paragraph{Lemma 4.}
Among strings with $\text{bal} < 0$, sorting by decreasing $\text{bal} - \text{minPref}$ is optimal.

\paragraph{Proof.}
Reverse each negative string and swap every parenthesis. This transforms it into a string with nonnegative
total balance. Under this transformation, the role of $\text{minPref}$ becomes exactly
\[
    \text{bal} - \text{minPref}.
\]
So by Lemma 3, the transformed strings should be sorted by decreasing that quantity. Reversing the
transformation gives exactly the claimed order for the original negative strings. \qed

\paragraph{Theorem.}
The algorithm outputs a valid permutation if and only if one exists.

\paragraph{Proof.}
By Lemma 2, it is enough to consider orders with all positive strings first and all negative strings last.
By Lemma 3, the chosen order is optimal for the positive group, and by Lemma 4 it is optimal for the
negative group. During the final scan, Lemma 1 exactly characterizes whether each next string can be
appended safely.

Therefore, if the algorithm fails, no valid ordering exists. If it succeeds, every prefix stays
nonnegative and the total balance is zero, so the produced order is a correct bracket sequence. \qed

\section*{Complexity Analysis}

Let $L$ be the total number of parentheses in the input. Computing all statistics takes $O(L)$ time.
Sorting the tablets takes $O(n \log n)$ time. The memory usage is $O(n)$.

\section*{Implementation Notes}

\begin{itemize}[leftmargin=*]
    \item We never need to store the full strings after computing $(\text{bal}, \text{minPref})$.
    \item The final total balance must be zero; otherwise a correct bracket sequence is impossible no matter
    how we permute the tablets.
\end{itemize}

\end{document}
